\documentclass[11pt,oneside]{article}
\usepackage[T1]{fontenc}
\usepackage[utf8]{inputenc}
\usepackage{lmodern}
\usepackage{microtype}
\usepackage[margin=1in]{geometry}
\usepackage{amsmath,amssymb,amsthm,mathtools}
\usepackage{booktabs,longtable,array,tabularx}
\usepackage{enumitem}
\usepackage{xcolor}
\usepackage{listings}
\usepackage{fancyhdr}
\usepackage[hypertexnames=false]{hyperref}
\usepackage{bookmark}
\usepackage{url}
\usepackage{graphicx}
\usepackage{pdfpages}
\usepackage{needspace}

\hypersetup{
  colorlinks=true,
  linkcolor=blue!55!black,
  citecolor=blue!55!black,
  urlcolor=blue!65!black,
  pdftitle={Radical Formulae Through Degree Four and Abel--Ruffini in iPRA},
  pdfauthor={OpenAI, prepared for Vladimir Reshetnikov}
}
\pagestyle{fancy}
\fancyhf{}
\lhead{Radical formulae and Abel--Ruffini in iPRA}
\rhead{\thepage}
\setlength{\headheight}{14pt}
\setlist{nosep,leftmargin=*}
\allowdisplaybreaks
\setlength{\emergencystretch}{3em}

\newtheorem{theorem}{Theorem}[section]
\newtheorem{lemma}[theorem]{Lemma}
\newtheorem{proposition}[theorem]{Proposition}
\newtheorem{corollary}[theorem]{Corollary}
\newtheorem{definition}[theorem]{Definition}
\newtheorem{remark}[theorem]{Remark}
\newtheorem{claim}[theorem]{Claim}

\newcommand{\N}{\mathbb N}
\newcommand{\Z}{\mathbb Z}
\newcommand{\Q}{\mathbb Q}
\newcommand{\C}{\mathbb C}
\newcommand{\F}{\mathbb F}
\newcommand{\Gal}{\operatorname{Gal}}
\newcommand{\Aut}{\operatorname{Aut}}
\newcommand{\sgn}{\operatorname{sgn}}
\newcommand{\rad}{\operatorname{Rad}}
\DeclareRobustCommand{\code}[1]{\ifmmode\text{\ttfamily\detokenize{#1}}\else\nolinkurl{#1}\fi}
\newcommand{\ipra}{\mathrm{iPRA}}
\newcommand{\poly}{\operatorname{Poly}}
\newcommand{\RatFun}{\operatorname{RatFun}}
\newcommand{\Der}{\operatorname{Der}}
\newcommand{\Comm}{\operatorname{Comm}}
\newcommand{\angles}[1]{\langle #1\rangle}

\lstdefinestyle{lean}{
  basicstyle=\ttfamily\small,
  columns=fullflexible,
  breaklines=true,
  frame=single,
  showstringspaces=false,
  keywordstyle=\bfseries,
  commentstyle=\itshape\color{black!65},
  keepspaces=true
}

\begin{document}

\begin{titlepage}
\centering
\vspace*{1.2in}
{\LARGE\bfseries Radical Formulae Through Degree Four\\[0.35em]
and the Abel--Ruffini Obstruction\\[0.35em]
in Intuitionistic Primitive Recursive Arithmetic\par}
\vspace{1.0in}
{\large A constructive, certificate-oriented proof and Lean formalization plan\par}
\vspace{0.7in}
{\large Prepared for Vladimir Reshetnikov\par}
\vfill
{\large 15 November 2026\par}
\vspace{0.4in}
\begin{minipage}{0.86\textwidth}
\small
This report is pinned to the ProveIt repository state inspected at commit
\code{25da87848ad4042c318af9bbc7f2e1c24614d6c3}.  Existing Lean files already
kernel-check the field-generic linear, quadratic, cubic, and quartic identities
and a semantic Abel--Ruffini obstruction.  The new contribution here is a
single constructive proof architecture showing how those results can be
expressed and verified inside the ProveIt presentation of intuitionistic
primitive recursive arithmetic, together with exact module boundaries for an
end-to-end Lean formalization.
\end{minipage}
\end{titlepage}
\pagenumbering{roman}

\begin{abstract}
We give exact formulae for polynomial equations of degrees at most four in a
form compatible with intuitionistic logic: a radical is never an unspecified
analytic branch, but an element accompanied by the equation it satisfies.  The
quadratic, Cardano, and Ferrari constructions are proved by complete
factorization identities, including all degenerate cases that normally hide
behind divisions by a radical.  We then prove the Abel--Ruffini theorem using
the generic quintic
\[
  \prod_{i=1}^{5}(T-x_i)
\]
over the subfield generated by its elementary symmetric coefficients.  Its
splitting-field Galois group is constructively identified with $S_5$.
Radical towers have solvable normal-closure groups, whereas an explicit
commutator calculation gives
$[S_5,S_5]=A_5$ and $[A_5,A_5]=A_5$; hence $S_5$ is not solvable.  It follows
that no complete formula by radicals exists in degree five, and padding by a
power of $T$ gives every degree at least five.

The entire argument is reduced to natural-number codes for finite sequences,
sparse polynomials, rational functions, permutations, finite group tables,
radical-tower records, and finite derivations.  Every operation and every
search used in the proof is primitive recursive, and every induction invariant
is a decidable equation, so quantifier-free induction suffices.  We state the
arithmetized theorems, explain their transfer through ProveIt's
\code{native_iff_pra} bi-interpretation, and give a detailed Lean file plan.
The archive also contains an exact symbolic verification script and its JSON
certificate report.
\end{abstract}

\tableofcontents
\clearpage
\pagenumbering{arabic}

\section{Statement, scope, and constructive reading}

\subsection{The two main results}

The phrase ``formula by radicals'' is often stated informally.  In a
constructive setting one must say what the radical symbols mean, what data are
being supplied, and what exactly is proved.  The following definition is the
one used throughout.

\begin{definition}[Radical tower and complete radical factorization]
Let $F$ be a field of characteristic zero.  A \emph{radical tower over $F$}
is a finite chain
\[
 F=F_0\subseteq F_1\subseteq\cdots\subseteq F_r
\]
with elements $\alpha_i\in F_i$ and integers $m_i\ge 2$ such that
\[
 F_i=F_{i-1}(\alpha_i),\qquad \alpha_i^{m_i}\in F_{i-1}.
\]
A polynomial $f(T)=a_nT^n+\cdots+a_0\in F[T]$, $a_n\ne0$, has a
\emph{complete radical factorization} if some radical tower contains elements
$r_1,\ldots,r_n$ satisfying the exact identity
\[
 f(T)=a_n\prod_{j=1}^{n}(T-r_j).
\]
Multiplicity is retained.  No ordering or analytic branch convention is part
of the definition.
\end{definition}

The identity is stronger than the assertion that each displayed value happens
to be a root: it gives soundness, exhaustiveness, and multiplicity in one
finite certificate.

\begin{theorem}[Constructive formula theorem through degree four]
\label{thm:low-degree-main}
For every characteristic-zero field and every nonzero polynomial of degree at
most four, the classical linear, quadratic, Cardano, and Ferrari constructions
admit the following intuitionistically valid formulation.

Given explicit radical witnesses satisfying the equations displayed in
Sections~\ref{sec:quadratic}, \ref{sec:cubic}, and \ref{sec:quartic}, a fixed
primitive-recursive expression returns a list of the required length, and an
exact factorization identity proves that this list contains precisely all
roots with multiplicity.  Compatible radical witnesses can be organized in a
radical tower.  No proof uses excluded middle for an arbitrary proposition;
the only case distinctions are equality tests on coded field expressions or
natural numbers.
\end{theorem}

\begin{theorem}[Constructive Abel--Ruffini theorem]
\label{thm:abel-main}
There is no complete formula by radicals for the general polynomial of degree
$n$ when $n\ge5$.  More precisely, let $x_1,\ldots,x_5$ be independent
indeterminates, let
\[
 L=\Q(x_1,\ldots,x_5),
\]
let $e_i$ be the elementary symmetric functions, and let
$K=\Q(e_1,\ldots,e_5)\subseteq L$.  Then
\[
 P(T)=\prod_{i=1}^{5}(T-x_i)\in K[T]
\]
has no complete radical factorization over $K$.  Consequently a coefficient-
parametrized radical solver valid for every quintic cannot exist.  For every
$n>5$, the polynomial $T^{n-5}P(T)$ gives the same obstruction in exact degree
$n$.

After coding all finite objects by natural numbers, the proof is formalizable
in intuitionistic primitive recursive arithmetic.  In particular, ProveIt's
native concise calculus proves the translation of the arithmetized theorem by
the established bi-interpretation with iPRA.
\end{theorem}

\subsection{What Abel--Ruffini does and does not say}

The negative theorem concerns a \emph{uniform complete} radical solution for
all polynomials of a fixed degree.  It does not say that every individual
quintic is unsolvable.  For example, $T^5-1$ is already solved by roots of
unity, and a reducible quintic may reduce to equations of smaller degree.  It
also does not forbid algorithms using operations beyond radicals, nor does it
forbid numerical approximation.

The word \emph{complete} matters.  A polynomial can have one easy root and
other roots that are not obtainable by radicals.  A purported solution of the
equation must account for the whole factorization, not merely produce one
root.  The generic polynomial used here is stronger still: its Galois group is
$S_5$, and no radical tower can contain all five independent roots.

\subsection{Why radical witnesses replace branch functions}

A complex square root is not a single-valued analytic function on all of
$\C$.  A cube root has three values, and Cardano's formula fails if two cube
roots are selected independently.  Therefore the constructive statement is
relational:
\[
  s^2=D,
  \qquad u^3=R,
  \qquad v^3=S,
  \qquad uv=-p/3.
\]
The data are the elements $s,u,v$ together with proofs of these equations.
The formula then proves a polynomial identity.  This is exactly the style of
the existing ProveIt Lean development: square and cube roots are represented
by field elements plus their defining equations, and the formal theorems are
generic over characteristic-zero fields.

This witness-oriented formulation is more constructive than choosing a global
branch.  It is also stronger algebraically: it works in every field extension
where the witnesses exist, not only in $\C$ and not only for a preferred
analytic branch.

\subsection{Relationship to the current ProveIt development}

At the inspected repository revision, the following semantic Lean layers are
already present.

\begin{itemize}
\item \code{Algebra/PolynomialFormulas/Lean/PolynomialFormulas/Basic.lean}
  proves the linear and quadratic formulae, exact quadratic factorization, and
  quadratic exhaustiveness.
\item \code{Cubic.lean} proves the depressed-cubic transformation, compatible
  Cardano radicals, the three-factor identity, and
  \code{cubic_eq_zero_iff}.
\item \code{Quartic.lean} proves Ferrari's resolvent equations, the two-
  quadratic factorization, and \code{quartic_eq_zero_iff}.
\item \code{GenericAbelRuffini.lean} constructs a generic polynomial whose
  Galois group is a symmetric group and derives a rootwise radical
  obstruction through mathlib's field-theoretic infrastructure.
\item \code{Logic/PrimitiveRecursiveArithmetic/ConciseIPRA} supplies the
  formal iPRA presentation and the theorem \code{native_iff_pra} transferring
  derivability between the native calculus and ordinary iPRA.
\end{itemize}

The present report does not replace those files.  It reorganizes their
mathematics into a deliberately finite, certificate-oriented proof whose
arithmetization uses only primitive recursion and quantifier-free induction.
Section~\ref{sec:lean-plan} identifies the exact additional bridge modules
needed for a fully kernel-checked iPRA result.

\subsection{Constructive discipline used below}

The proof observes four rules.

\begin{enumerate}
\item A radical is always accompanied by its power equation.
\item Division is performed only with an explicit nonzero hypothesis.
\item A disjunction is introduced only after a decidable finite equality test,
  or it is avoided by retaining a factorization identity.
\item Every finite construction is given with a primitive-recursive bound;
  no appeal is made to unbounded minimization, a global choice function, or a
  completed infinite set.
\end{enumerate}

These rules will be repeated at the few places where a classical textbook
proof usually suppresses a branch or an existence argument.

\section{The iPRA framework}
\label{sec:ipra-framework}

\subsection{Why an arithmetic theory can discuss fields and Galois groups}

The only primitive objects of iPRA are natural numbers.  Algebraic objects are
therefore represented by finite codes.  This is not a limitation for the
present proof.  Every polynomial, rational function, permutation, group table,
radical tower, field expression, and derivation appearing below is finite.
The semantic notation is a readable abbreviation for operations on such
codes.

The ProveIt presentation of iPRA has primitive-recursive function symbols and
quantifier-free induction.  Its native scalar calculus is bi-interpretable
with the ordinary first-order presentation; the numerical domain remains all
of $\N$ and equality remains actual equality.  Consequently it is enough to
show the following two facts.

\begin{enumerate}
\item Each construction used in the proof is primitive recursive on codes.
\item Each induction can be stated as an equation or a Boolean invariant about
  a primitive-recursive computation.
\end{enumerate}

The proof does not need induction on an arbitrary semantic property of fields.
It needs induction on lengths of lists, degrees of sparse polynomials, sizes of
finite groups, and lengths of derivations.  All these invariants are decidable.

\subsection{A concrete coding layer}

One convenient coding, not the only possible one, is as follows.

\begin{longtable}{p{0.22\textwidth}p{0.70\textwidth}}
\toprule
Object & Natural-number representation and primitive-recursive operations\\
\midrule
\endhead
Finite sequence & Iterated pairing with stored length.  Entry, append, map,
fold, bounded search, and permutation are primitive recursive.\\
Integer & Zig-zag code: $0,1,-1,2,-2,\ldots$.  Addition and multiplication are
primitive recursive.\\
Rational number & A signed numerator and positive denominator, optionally
normalized by Euclid's algorithm.  Equality is cross multiplication.\\
Monomial & A coefficient and a finite exponent vector.\\
Sparse polynomial & A sorted list of nonzero monomials.  Normalization combines
equal exponent vectors; addition, multiplication, substitution, and formal
derivative are primitive recursive.\\
Rational function & A pair $(P,Q)$ with $Q\ne0$.  Equality of $(P,Q)$ and
$(R,S)$ is the polynomial equation $PS=RQ$.\\
Permutation of $d$ letters & A length-$d$ list containing each number below
$d$ exactly once.  Composition, inverse, sign, cycle decomposition, and
bounded enumeration of all $d!$ permutations are primitive recursive.\\
Finite group & A finite carrier list with multiplication, inverse, and identity
tables, or a subgroup of a permutation group represented by a Boolean
membership table.\\
Field expression & A syntax tree containing constants, generators, $+,-,\cdot$,
inversion with a nonzero side certificate, and named algebraic generators.\\
Radical step & An exponent $m\ge2$, a radicand expression $a$, a new generator
$z$, and the relation $z^m=a$.\\
Derivation & A list of proof nodes in a fixed finite calculus; every node names
its premises and has a local primitive-recursive checker.\\
\bottomrule
\end{longtable}

There is no need to decide all true identities in an arbitrary finitely
presented field.  A certificate carries an explicit derivation.  The checker
only verifies local rules: normalized ring identities, substitution of a
named relation, use of an explicit inverse equation, finite linear algebra,
and finite group computations.  Soundness of the checker is proved by
quantifier-free induction on the derivation length.

\subsection{Proof-carrying field presentations}

A \emph{proof-carrying presentation} consists of:

\begin{itemize}
\item a finite list of generators;
\item a finite list of polynomial and inverse relations;
\item syntax trees for every element used in the argument;
\item a derivation for every asserted equality or nonzero condition;
\item for a finite extension, a finite basis and coordinate calculations for
  multiplication, embeddings, and automorphisms used by the proof.
\end{itemize}

The checker is primitive recursive because it never searches for a proof.
It checks the supplied proof.  This is the same design principle used by a
small proof assistant: semantic truth may be difficult, but checking a finite
proof object is mechanical.

For the generic quintic the base field has an especially simple direct
representation.  Start with sparse polynomials in five variables over $\Q$,
then represent $L=\Q(x_1,\ldots,x_5)$ by fractions.  Equality in $L$ is a
single cross-multiplied polynomial identity.  The subfield
$K=\Q(e_1,\ldots,e_5)$ is represented by rational expressions in five named
symmetric generators together with their embeddings in $L$.  All operations
used in the proof are therefore explicit normalization computations.

\subsection{Radical programs as relational data}

A global operation $z\mapsto\sqrt[m]{z}$ is inappropriate constructively.
Instead we use a finite sequential program with a nondeterministic but
proof-carrying radical instruction:
\[
  z_i^m = a_i(z_0,\ldots,z_{i-1}).
\]
A \emph{run} is a tuple of values in an extension field satisfying every
instruction.  Arithmetic instructions are deterministic; radical
instructions supply a value and a proof of its power equation.  A run is
\emph{complete for $f$} if the designated outputs give the exact factorization
of $f$.

Every ordinary nested-radical expression compiles to such a program.  A case
distinction compiles to a finite decision tree whose guards are decidable
equalities of normalized expressions.  Conversely, a run generates a radical
tower by adjoining its radical values in program order.  Thus the relational
language is neither weaker nor stronger than the standard field-theoretic
notion needed for Abel--Ruffini; it merely exposes the witnesses.

\subsection{Arithmetic predicates used in the final theorem}

The formalization can use the following coded predicates.

\begin{align*}
 \operatorname{RatFunField}(k) &:\quad k\text{ codes the base rational-function presentation},\\
 \operatorname{RadTower}(k,t) &:\quad t\text{ codes a valid radical tower over }k,\\
 \operatorname{Splits}(k,p,t) &:\quad p\text{ factors completely in the top of }t,\\
 \operatorname{Formula}(n,c) &:\quad c\text{ is a degree-}n\text{ radical-program code},\\
 \operatorname{Run}(k,c,w) &:\quad w\text{ is a valid run of }c\text{ over }k.
\end{align*}

The local checking relations can be made primitive recursive.  Their semantic
soundness is a theorem of iPRA about the checker.  The negative theorem is
conveniently factored into two statements:

\begin{align}
 &\forall t\;\neg\operatorname{Splits}(K,P,t),
 \label{eq:no-tower}\\
 &\operatorname{Formula}(5,c)\land\operatorname{Universal}(c)
 \;\longrightarrow\;
   \exists t\;\operatorname{Splits}(K,P,t).
 \label{eq:compile-formula}
\end{align}

Equation~\eqref{eq:no-tower} is the Galois obstruction.
Equation~\eqref{eq:compile-formula} is the primitive-recursive compiler that
specializes a purported coefficient formula to the generic coefficients.
Together they imply that no universal code $c$ exists.

For degrees at most four, fixed constructors produce a radical program and a
factorization derivation from the supplied witness equations.  The internal
statements therefore have the form
\[
 \forall a_0,\ldots,a_n,w\;
   \bigl(\operatorname{WitnessEquations}_n(a,w)=1\bigr)
   \to
   \bigl(\operatorname{FactorCheck}_n(a,w)=1\bigr).
\]
Both sides are equations between primitive-recursive terms, exactly the kind
of induction-free or quantifier-free-induction statement natural to PRA.

\subsection{Where intuitionistic logic matters}

Three classical shortcuts are explicitly removed.

\paragraph{No arbitrary branch choice.}
The square or cube root is a witness, not the result of a choice function.
Cardano's two cube radicals are linked by a product equation.

\paragraph{No unrestricted zero/nonzero split.}
The human proof sometimes writes ``if $s=0$''.  Internally $s$ is a normalized
coded expression in a decidable field presentation, so the split is a Boolean
calculation.  When no decidable equality is available, the primary theorem is
kept as a factorization identity and makes no disjunctive claim.

\paragraph{No appeal to an infinite set of all automorphisms.}
The only automorphism groups used are finite.  For the generic quintic every
automorphism is represented by one of the $120$ permutations of five
symbols.  Radical normal-closure groups are represented by finite tables whose
sizes are bounded by products and factorials extracted from the tower code.

\subsection{Quantifier-free induction is enough}

Examples of the induction statements actually needed are:

\begin{itemize}
\item after processing the first $j$ monomials, the accumulator equals their
  normalized sum;
\item after processing the first $j$ proof nodes, each checked node is sound;
\item after processing the first $j$ radical steps, the generated group-series
  certificate passes its finite checker;
\item after $j$ closure rounds, the Boolean table equals the subgroup generated
  by commutators found so far.
\end{itemize}

Each invariant is an equation between values of primitive-recursive functions.
No induction on a formula containing an unbounded existential quantifier is
required.  Existential mathematical statements are represented by explicit
witness codes produced before the checker is invoked.

\section{Linear and quadratic equations}

\subsection{The linear equation}

Let
\[
 f(X)=aX+b,
 \qquad a\ne0.
\]
Define
\[
 r=-\frac{b}{a}.
\]
Then
\[
 f(r)=a\left(-\frac ba\right)+b=0,
 \qquad
 f(X)=a(X-r).
\]
The factorization proves both correctness and uniqueness.  The only division
uses the explicit hypothesis $a\ne0$.  As a coded computation, field inversion
is supplied with a proof that $a$ is nonzero; normalization verifies the two
equalities above.

\subsection{The quadratic formula}
\label{sec:quadratic}

Let
\[
 f(X)=aX^2+bX+c,
 \qquad a\ne0,
\]
and let $s$ be a supplied square-root witness for the discriminant:
\[
 s^2=b^2-4ac.
 \tag{Q1}\label{eq:quad-disc}
\]
Because the characteristic is zero, $2a\ne0$.  Define
\[
 r_+=\frac{-b+s}{2a},
 \qquad
 r_-=\frac{-b-s}{2a}.
 \tag{Q2}\label{eq:quad-roots}
\]

\begin{theorem}[Exact quadratic factorization]
\label{thm:quadratic-factorization}
Under $a\ne0$ and~\eqref{eq:quad-disc},
\[
 aX^2+bX+c=a(X-r_+)(X-r_-).
 \tag{Q3}\label{eq:quad-factor}
\]
Consequently $r_+$ and $r_-$ are roots, and in a coded field with decidable
equality every root equals one of them.
\end{theorem}

\begin{proof}
The two elementary symmetric functions of the displayed values are
\[
 r_++r_-=
 \frac{-2b}{2a}=-\frac ba.
\]
and
\[
 r_+r_-=
 \frac{(-b+s)(-b-s)}{4a^2}
 =\frac{b^2-s^2}{4a^2}
 =\frac{b^2-(b^2-4ac)}{4a^2}
 =\frac ca.
\]
Therefore
\[
 a(X-r_+)(X-r_-)
 =aX^2-a(r_++r_-)X+a r_+r_-
 =aX^2+bX+c.
\]
This is a polynomial identity derived only from~\eqref{eq:quad-disc}.
If $f(x)=0$, then
$a(x-r_+)(x-r_-)=0$.  Since $a\ne0$, one has
$(x-r_+)(x-r_-)=0$.  In the coded field, equality with zero is decidable:
if $x-r_+=0$ we have $x=r_+$; otherwise its inverse exists and multiplication
by that inverse gives $x-r_-=0$.  Thus $x=r_-$.
\end{proof}

The proof remains valid when $s=0$; then the two list entries coincide and the
factorization records the double root.  No separate discriminant case is
needed.

\subsection{Primitive-recursive certificate}

The quadratic certificate has the following fixed shape.

\begin{enumerate}
\item Parse codes for $a,b,c,s$.
\item Check the nonzero certificate for $a$.
\item Normalize the two sides of $s^2=b^2-4ac$.
\item Construct syntax trees for $r_+$ and $r_-$.
\item Clear the denominator $(2a)^2$ in~\eqref{eq:quad-factor}.
\item Normalize the resulting integer-coefficient polynomial expression and
  reduce once by the relation~\eqref{eq:quad-disc}.
\end{enumerate}

Every step is primitive recursive.  The accompanying Python verifier performs
exactly this calculation and obtains remainder zero.  In Lean, ProveIt's
\code{quadratic_formula_factorization} theorem already provides the semantic
field-generic result, and \code{quadratic_eq_zero_iff} gives exhaustiveness.

\section{Cardano's cubic formula}
\label{sec:cubic}

\subsection{Depressing the cubic}

Let
\[
 f(X)=aX^3+bX^2+cX+d,
 \qquad a\ne0.
\]
Put
\[
 A=\frac ba,\qquad B=\frac ca,\qquad C=\frac da,
\]
so that $f(X)=a(X^3+AX^2+BX+C)$.  Make the translation
\[
 X=Y-\frac A3.
\]
Direct expansion gives
\[
 X^3+AX^2+BX+C=Y^3+pY+q,
 \tag{C1}\label{eq:cubic-depress}
\]
where
\[
 p=B-\frac{A^2}{3}
   =\frac{3ac-b^2}{3a^2},
 \qquad
 q=\frac{2A^3}{27}-\frac{AB}{3}+C
   =\frac{2b^3-9abc+27a^2d}{27a^3}.
 \tag{C2}\label{eq:cubic-pq}
\]
The denominators are nonzero in characteristic zero.  It remains to factor the
depressed cubic
\[
 g(Y)=Y^3+pY+q.
\]

\subsection{The compatible radical equations}

Choose $u,v$ satisfying
\[
 u^3+v^3=-q,
 \qquad
 uv=-\frac p3.
 \tag{C3}\label{eq:cardano-compatible}
\]
Then
\[
 (u+v)^3+p(u+v)+q
 =u^3+v^3+3uv(u+v)+p(u+v)+q=0.
\]
Thus $u+v$ is a root.  The product equation in
\eqref{eq:cardano-compatible} is essential.  Independent choices of cube-root
branches do not in general satisfy it.

Let $\omega$ be a primitive cube-root witness in the algebraic form
\[
 \omega^2+\omega+1=0.
 \tag{C4}\label{eq:omega}
\]
Then $\omega^3=1$ and $1+\omega+\omega^2=0$.  In characteristic zero,
$\omega\ne1$.  Such a witness itself uses only a square radical: if
$\tau^2=-3$, then $\omega=(-1+\tau)/2$ satisfies~\eqref{eq:omega}.

Define
\[
 r_0=u+v,
 \quad
 r_1=\omega u+\omega^2v,
 \quad
 r_2=\omega^2u+\omega v.
 \tag{C5}\label{eq:cardano-roots-depressed}
\]

\begin{theorem}[Cardano factorization]
\label{thm:cardano-factorization}
From~\eqref{eq:cardano-compatible} and~\eqref{eq:omega} one can derive
\[
 Y^3+pY+q=(Y-r_0)(Y-r_1)(Y-r_2).
 \tag{C6}\label{eq:cardano-factorization}
\]
Consequently the roots of the original cubic are
\[
 x_k=r_k-\frac A3,
 \qquad k=0,1,2.
 \tag{C7}\label{eq:cardano-original-roots}
\]
\end{theorem}

\begin{proof}
The sum of the three displayed values is
\[
 r_0+r_1+r_2
 =(1+\omega+\omega^2)u+(1+\omega+\omega^2)v=0.
\]
A direct expansion, reduced by $\omega^2+\omega+1=0$, gives
\[
 r_0r_1+r_0r_2+r_1r_2=-3uv=p.
\]
Similarly,
\[
 r_0r_1r_2=u^3+v^3=-q.
\]
The monic polynomial with these three roots is therefore
\[
 Y^3-(r_0+r_1+r_2)Y^2
 +(r_0r_1+r_0r_2+r_1r_2)Y-r_0r_1r_2
 =Y^3+pY+q.
\]
Combining this identity with~\eqref{eq:cubic-depress} proves the claim.
\end{proof}

The factorization, rather than three independent substitution calculations,
is the preferred certificate.  It proves all roots simultaneously and retains
multiplicity when two or three values coincide.

\subsection{Constructing a compatible pair without a hidden pole}

The usual displayed formula introduces
\[
 \Delta=\left(\frac q2\right)^2+\left(\frac p3\right)^3
 \tag{C8}\label{eq:cardano-delta}
\]
and a square-root witness $s^2=\Delta$.  Put
\[
 R=-\frac q2+s,
 \qquad
 S=-\frac q2-s.
 \tag{C9}\label{eq:RS}
\]
Then
\[
 R+S=-q,
 \qquad
 RS=\frac{q^2}{4}-s^2=-\left(\frac p3\right)^3.
 \tag{C10}\label{eq:RS-relations}
\]
A textbook often says ``take $u^3=R$ and $v^3=S$'' without controlling the
branches.  The following finite construction is correct in every case.

\paragraph{Case $p\ne0$.}
Equation~\eqref{eq:RS-relations} implies $R\ne0$ and $S\ne0$.  Adjoin a cube
root $u$ with $u^3=R$ and define
\[
 v=-\frac{p}{3u}.
 \tag{C11}\label{eq:v-defined}
\]
The denominator is nonzero.  Then $uv=-p/3$, and
\[
 v^3=\frac{(-p/3)^3}{u^3}
 =\frac{-(p/3)^3}{R}=S
\]
by $RS=-(p/3)^3$.  Hence $u^3+v^3=R+S=-q$.

\paragraph{Case $p=0$.}
Now $RS=0$.  Adjoin cube-root witnesses independently,
\[
 u^3=R,\qquad v^3=S.
\]
Then
\[
 (uv)^3=u^3v^3=RS=0.
\]
A field has no nonzero nilpotents, so $uv=0=-p/3$; also
$u^3+v^3=R+S=-q$.  No decision whether $R$ or $S$ vanishes is needed.

Thus the construction uses one square radical, at most two cube radicals, and
the square radical used to construct $\omega$.  The $p\ne0$ branch needs only
one independently chosen cube root because the other is defined by division;
the $p=0$ branch uses two cube-root witnesses but contains no division by a
possibly zero quantity.

\subsection{The familiar displayed formula}

With a compatible choice of cube roots, the three roots may be written
\[
 x_k=-\frac{b}{3a}
 +\omega^k\sqrt[3]{-\frac q2+\sqrt\Delta}
 +\omega^{-k}\sqrt[3]{-\frac q2-\sqrt\Delta},
 \qquad k=0,1,2,
 \tag{C12}\label{eq:classical-cardano}
\]
where the two cube radicals are constrained so that their product is $-p/3$.
Equation~\eqref{eq:v-defined} is a branch-safe implementation of that
constraint when $p\ne0$; the $p=0$ branch above removes the apparent pole.

\subsection{Constructive exhaustiveness}

If $g(y)=0$, equation~\eqref{eq:cardano-factorization} gives
\[
 (y-r_0)(y-r_1)(y-r_2)=0.
\]
In a coded field, test $y-r_0=0$.  If false, cancel its explicit inverse and
test $y-r_1=0$; if that is false, cancellation gives $y-r_2=0$.  This is a
finite decision tree, not an application of excluded middle to an arbitrary
proposition.  The same argument after translation proves that every root of
$f$ occurs in~\eqref{eq:cardano-original-roots}.

\subsection{Primitive-recursive verification}

The certificate checker expands the right-hand side of
\eqref{eq:cardano-factorization} and reduces it by the three relations
\[
 \omega^2+\omega+1=0,
 \qquad
 3uv+p=0,
 \qquad
 u^3+v^3+q=0.
\]
Because the degrees and number of variables are fixed, even a naive sparse
normalizer is primitive recursive.  The supplied verification script computes
a Groebner normal form only as an independent convenience check; the formal
iPRA proof needs merely a fixed sequence of ring rewrites, exactly as the
existing Lean theorem \code{cardano_factorization} does.  ProveIt's
\code{cubic_eq_zero_iff} already packages the semantic soundness and
exhaustiveness statement over a characteristic-zero field.

\section{Ferrari's quartic formula}
\label{sec:quartic}

\subsection{Depressing the quartic}

Let
\[
 f(X)=aX^4+bX^3+cX^2+dX+e,
 \qquad a\ne0.
\]
Put
\[
 A=\frac ba,\quad B=\frac ca,\quad C=\frac da,\quad D=\frac ea
\]
and translate
\[
 X=Y-\frac A4.
\]
A direct expansion gives
\[
 \frac1a f\left(Y-\frac A4\right)=Y^4+pY^2+qY+r,
 \tag{F1}\label{eq:quartic-depress}
\]
where
\begin{align}
 p&=B-\frac{3A^2}{8},\label{eq:quartic-p}\\
 q&=C-\frac{AB}{2}+\frac{A^3}{8},\label{eq:quartic-q}\\
 r&=D-\frac{AC}{4}+\frac{A^2B}{16}-\frac{3A^4}{256}.
 \label{eq:quartic-r}
\end{align}
We therefore solve
\[
 h(Y)=Y^4+pY^2+qY+r.
\]

\subsection{The biquadratic branch}
\label{sec:quartic-biquadratic}

When $q=0$, no cubic resolvent is needed.  Apply the quadratic formula to
\[
 Z^2+pZ+r=0.
\]
Choose $\delta,z_+,z_-$ with
\[
 \delta^2=p^2-4r,
 \qquad
 z_+=\frac{-p+\delta}{2},
 \qquad
 z_-=\frac{-p-\delta}{2}.
\]
Then
\[
 Z^2+pZ+r=(Z-z_+)(Z-z_-).
\]
Choose square-root witnesses $\eta_+^2=z_+$ and $\eta_-^2=z_-$.  The exact
factorization is
\[
 Y^4+pY^2+r
 =(Y-\eta_+)(Y+\eta_+)(Y-\eta_-)(Y+\eta_-).
 \tag{F2}\label{eq:biquadratic-factor}
\]
This branch is entirely constructive once the equation $q=0$ is supplied; it
does not test equality of a newly introduced radical.

\subsection{Ferrari's factorization equations when \texorpdfstring{$q\ne0$}{q nonzero}}

Assume now that $q\ne0$.  Seek a factorization into two monic quadratics of
opposite linear terms:
\[
 h(Y)=\bigl(Y^2-sY+(m-t)\bigr)
       \bigl(Y^2+sY+(m+t)\bigr).
 \tag{F3}\label{eq:ferrari-target}
\]
Multiplication shows that~\eqref{eq:ferrari-target} follows from
\[
 s^2=2m-p,
 \qquad
 2st=-q,
 \qquad
 t^2=m^2-r.
 \tag{F4}\label{eq:ferrari-parameters}
\]
Indeed, the $Y^2$ coefficient is $2m-s^2=p$, the $Y$ coefficient is
$-2st=q$, and the constant term is $m^2-t^2=r$.

Eliminating $s$ and $t$ gives the cubic resolvent
\[
 q^2-4(2m-p)(m^2-r)=0.
 \tag{F5}\label{eq:ferrari-resolvent}
\]
Equivalently,
\[
 m^3-\frac p2m^2-rm+\frac{pr}{2}-\frac{q^2}{8}=0.
 \tag{F6}\label{eq:ferrari-resolvent-monic}
\]
Thus $m$ is obtained by the cubic construction of
Section~\ref{sec:cubic}.

Choose a resolvent root $m$ and a square-root witness
\[
 s^2=2m-p.
 \tag{F7}\label{eq:s-square}
\]
The hypothesis $q\ne0$ and~\eqref{eq:ferrari-resolvent} imply
$2m-p\ne0$: otherwise the right-hand product in
\eqref{eq:ferrari-resolvent} would vanish and hence $q^2=0$, forcing $q=0$.
Consequently $s\ne0$.  We may therefore define
\[
 t=-\frac{q}{2s}.
 \tag{F8}\label{eq:t-quotient}
\]
Then $2st=-q$.  The resolvent equation and~\eqref{eq:s-square} give
\[
 q^2=4(2m-p)(m^2-r)=4s^2(m^2-r),
\]
so
\[
 t^2=\frac{q^2}{4s^2}=m^2-r.
\]
All three equations in~\eqref{eq:ferrari-parameters} hold, and every division
has an explicit nonzero proof.

\subsection{Solving the two quadratics}

Choose square-root witnesses
\[
 \rho^2=s^2-4(m-t),
 \qquad
 \sigma^2=s^2-4(m+t).
 \tag{F9}\label{eq:rho-sigma}
\]
The first factor in~\eqref{eq:ferrari-target} has roots
\[
 y_0=\frac{s+\rho}{2},
 \qquad
 y_1=\frac{s-\rho}{2},
 \tag{F10}\label{eq:quartic-first-pair}
\]
and the second has roots
\[
 y_2=\frac{-s+\sigma}{2},
 \qquad
 y_3=\frac{-s-\sigma}{2}.
 \tag{F11}\label{eq:quartic-second-pair}
\]
The four roots of the original quartic are
\[
 x_i=y_i-\frac A4,
 \qquad i=0,1,2,3.
 \tag{F12}\label{eq:quartic-roots}
\]

\begin{theorem}[Exact Ferrari factorization]
\label{thm:ferrari-factorization}
Assume~\eqref{eq:ferrari-parameters} and~\eqref{eq:rho-sigma}.  Then
\[
 Y^4+pY^2+qY+r
 =\prod_{i=0}^{3}(Y-y_i),
 \tag{F13}\label{eq:quartic-complete-factor}
\]
with $y_i$ as in~\eqref{eq:quartic-first-pair} and
\eqref{eq:quartic-second-pair}.  After the translation~\eqref{eq:quartic-depress},
\[
 f(X)=a\prod_{i=0}^{3}(X-x_i).
 \tag{F14}\label{eq:quartic-original-factor}
\]
\end{theorem}

\begin{proof}
The equations~\eqref{eq:ferrari-parameters} give the two-quadratic identity
\eqref{eq:ferrari-target}.  By the quadratic factorization theorem,
\[
 Y^2-sY+(m-t)=(Y-y_0)(Y-y_1)
\]
because its discriminant is $s^2-4(m-t)=\rho^2$, and
\[
 Y^2+sY+(m+t)=(Y-y_2)(Y-y_3)
\]
because its discriminant is $s^2-4(m+t)=\sigma^2$.  Multiplying these
identities proves~\eqref{eq:quartic-complete-factor}.  Substituting
$Y=X+A/4$ and multiplying by $a$ gives
\eqref{eq:quartic-original-factor}.
\end{proof}

\subsection{Completeness and multiplicity}

In the $q=0$ branch, equation~\eqref{eq:biquadratic-factor} is already the
complete four-factor certificate.  In the $q\ne0$ branch, any root of the
depressed quartic makes the product of the two quadratic factors zero.  The
field no-zero-divisors law reduces the problem to one of the two quadratic
equations; Theorem~\ref{thm:quadratic-factorization} then yields one of the
four values.  Repeated roots cause no difficulty: the list may contain equal
entries, while the factorization retains their multiplicities.

\subsection{Primitive-recursive verification}

The certificate checker first verifies the translation
\eqref{eq:quartic-depress}.  It then reads the supplied branch certificate:
$q=0$ selects the biquadratic identity; $q\ne0$ selects Cardano's cubic
resolver and Ferrari's equations.  In the latter branch it checks
\[
 s^2-(2m-p)=0,
 \quad 2st+q=0,
 \quad t^2-(m^2-r)=0,
\]
then expands~\eqref{eq:ferrari-target} and applies the already verified
quadratic certificates.  This composition is primitive recursive.  The
existing Lean file \code{Quartic.lean} proves precisely the parametric Ferrari
identity and its soundness/exhaustiveness once $m,s,t,\rho,\sigma$ are
supplied; the additional biquadratic branch is a direct composition of its
quadratic layer.

\section{Why radical towers have solvable Galois groups}
\label{sec:radical-solvable}

This section proves the algebraic implication used by Abel--Ruffini.  The proof
is arranged so that every set under discussion is finite and every map can be
recorded by a finite table.  The usual field-theoretic language is retained
for readability; Section~\ref{sec:internalization} replaces each object by its
natural-number code.

\subsection{Finite-group solvability}

For a finite group $G$, let $G'=[G,G]$ be the subgroup generated by all
commutators
\[
 [g,h]=ghg^{-1}h^{-1}.
\]
Define the derived series by
\[
 G^{(0)}=G,
 \qquad
 G^{(i+1)}=(G^{(i)})'.
\]
The group is \emph{solvable} if $G^{(r)}=1$ for some $r$.
For a finite table this is decidable: compute commutators, close under
multiplication and inverse, and iterate.  The subgroup chain can strictly
decrease at most $|G|-1$ times, so $|G|$ iterations suffice.

We use three closure facts, each admitting a finite constructive proof.

\begin{lemma}[Subgroups]
A subgroup of a solvable finite group is solvable.
\end{lemma}

\begin{proof}
If $H\le G$, then $H'\le G'$ because every commutator of elements of $H$ is a
commutator in $G$.  Primitive recursion gives $H^{(i)}\le G^{(i)}$ for all
$i$.  Once $G^{(r)}=1$, also $H^{(r)}=1$.
\end{proof}

\begin{lemma}[Quotients]
A quotient of a solvable finite group is solvable.
\end{lemma}

\begin{proof}
For a surjective homomorphism $\pi:G\to Q$, one has
$\pi(G')=Q'$: images of commutators are commutators, and every commutator in
$Q$ has lifted arguments.  Hence $\pi(G^{(i)})=Q^{(i)}$ by finite induction.
If $G^{(r)}=1$, then $Q^{(r)}=1$.
\end{proof}

\begin{lemma}[Extensions]
Suppose
\[
 1\longrightarrow N\longrightarrow G\longrightarrow Q\longrightarrow1
\]
is exact and both $N$ and $Q$ are solvable.  Then $G$ is solvable.
\end{lemma}

\begin{proof}
If $Q^{(r)}=1$, the quotient calculation gives $G^{(r)}\le N$.  If
$N^{(s)}=1$, then
\[
 G^{(r+s)}=(G^{(r)})^{(s)}\le N^{(s)}=1.
\]
All subgroup inclusions can be verified from finite multiplication tables.
\end{proof}

\subsection{One normal radical layer}

Let $F$ be a characteristic-zero field containing all $m$th roots of unity,
and suppose $E/F$ is a finite normal extension generated by roots of finitely
many polynomials
\[
 Z^m-a_1,\ldots,Z^m-a_k,
 \qquad a_j\in F.
\]
Choose roots $\beta_j$ with $\beta_j^m=a_j$.  If $a_j\ne0$, every root of
$Z^m-a_j$ is $\zeta\beta_j$ for some $m$th root of unity $\zeta\in F$.
If $a_j=0$, the only root is zero and that generator is redundant.

For $\sigma\in\Gal(E/F)$, each nonzero $\beta_j$ is therefore sent to
\[
 \sigma(\beta_j)=\zeta_j(\sigma)\beta_j,
 \qquad \zeta_j(\sigma)^m=1.
\]
The map
\[
 \sigma\longmapsto
 (\zeta_1(\sigma),\ldots,\zeta_k(\sigma))
\]
is an injective group homomorphism into a finite product of cyclic groups of
$m$th roots of unity.  It is injective because the $\beta_j$ generate $E$ over
$F$.  Hence $\Gal(E/F)$ is abelian.

This argument is finite.  A certificate lists the generators, their power
equations, and the root-of-unity multiplier by which each automorphism sends
each generator.  Homomorphism and injectivity are bounded table checks.

\subsection{A normal closure built one radical layer at a time}

Consider a radical tower
\[
 K=K_0\subseteq K_1\subseteq\cdots\subseteq K_r=M,
 \qquad
 K_i=K_{i-1}(\alpha_i),
 \qquad
 \alpha_i^{m_i}=a_i\in K_{i-1}.
 \tag{R1}\label{eq:radical-tower}
\]
Let $m$ be a common multiple of all $m_i$ and put
\[
 N_0=K(\mu_m),
\]
where $\mu_m$ denotes the $m$th roots of unity.  The extension $N_0/K$ is
finite Galois.  An automorphism is determined by the image of a primitive
$m$th root, so its Galois group embeds in the abelian group
$(\Z/m\Z)^\times$ and is therefore abelian.

We now construct finite extensions
\[
 N_0\subseteq N_1\subseteq\cdots\subseteq N_r=N
\]
such that every $N_i/K$ is normal and $N_i$ contains the images of
$N_0K_i$ under all $K$-embeddings into a fixed finite splitting presentation.
Assume $N_{i-1}$ has been constructed.  Let $A_i$ be the finite set
\[
 A_i=\bigl\{\tau(a_i):
   \tau:N_0K_{i-1}\hookrightarrow N_{i-1}
   \text{ is a $K$-embedding}\bigr\}.
\]
For each $b\in A_i$, adjoin one root $\beta_b$ of $Z^{m_i}-b$ and put
\[
 N_i=N_{i-1}(\beta_b:b\in A_i).
\]
Because $N_{i-1}$ contains all $m_i$th roots of unity, every polynomial
$Z^{m_i}-b$ splits in $N_i$.  The set $A_i$ is stable under
$K$-automorphisms of $N_{i-1}$: composing a listed embedding with such an
automorphism merely produces another listed embedding.  Consequently every
$K$-embedding of $N_i$ maps each adjoined root to another root already in
$N_i$.  Thus $N_i/K$ is normal.

The restriction kernel
\[
 H_i=\Gal(N_i/N_{i-1})
\]
is abelian by the preceding subsection: $N_i$ is generated over $N_{i-1}$ by
roots of the finitely many binomials $Z^{m_i}-b$, and $N_{i-1}$ contains the
required roots of unity.  Since both $N_i/K$ and $N_{i-1}/K$ are finite
normal separable extensions, restriction gives an exact sequence
\[
 1\longrightarrow H_i
 \longrightarrow \Gal(N_i/K)
 \longrightarrow \Gal(N_{i-1}/K)
 \longrightarrow1.
 \tag{R2}\label{eq:radical-exact}
\]
Surjectivity is the finite extension-of-embeddings theorem in this
special situation.  Regard an automorphism of $N_{i-1}$ as a $K$-embedding
into $N_i$ and extend it across the finitely generated algebraic extension
$N_i/N_{i-1}$.  Because $N_i/K$ is normal, the extended embedding has image
$N_i$ and is therefore an automorphism.  One must not choose the images of the
$\beta_b$ independently, since they may satisfy relations.  In the coded
proof, a single coordinate map for the extension is supplied; the checker
verifies all generator relations, multiplicativity, and invertibility.  Thus
no choice principle is invoked by the verifier.

Starting with the abelian group $\Gal(N_0/K)$ and repeatedly applying the
extension lemma to~\eqref{eq:radical-exact}, we obtain that
$\Gal(N/K)$ is solvable.  The field $N$ contains the original tower $M$ via
the chosen embedding of each $\alpha_i$ into the corresponding root layer.

\begin{theorem}[Radical-tower theorem]
\label{thm:radical-group-solvable}
Let $M/K$ be a finite extension contained in a finite radical tower of
characteristic-zero fields.  There is a finite normal extension $N/K$
containing $M$ such that $\Gal(N/K)$ is solvable.
\end{theorem}

\begin{proof}
Remove any trivial zero radical generators, adjoin the finitely many required
roots of unity, and apply the layer construction above.  Characteristic zero
gives separability.  The exact sequences have abelian kernels and the initial
cyclotomic group is abelian, so finite induction proves solvability.
\end{proof}

\subsection{Passing to a splitting field}

\begin{corollary}[A polynomial split by radicals has solvable Galois group]
\label{cor:split-rad-solvable}
Let $f\in K[T]$ be separable, and let $L/K$ be its splitting field.  If a
radical tower contains all roots of $f$, then $\Gal(L/K)$ is solvable.
\end{corollary}

\begin{proof}
Take the solvable normal extension $N/K$ from
Theorem~\ref{thm:radical-group-solvable}; it contains $L$.  Since $L/K$ is
normal, restriction gives a surjection
\[
 \Gal(N/K)\twoheadrightarrow\Gal(L/K).
\]
The target is a quotient of a solvable group and is therefore solvable.
\end{proof}

The extension steps are constructive in the certificate sense used here: all
extensions are finite, every relevant polynomial has a finite explicitly
listed root set, and every embedding is extended one generator at a time by a
supplied finite table.  The checker uses neither a maximal-ideal theorem nor a
global algebraic closure.

\subsection{Primitive-recursive bounds}

Let the tower length be $r$ and let the exponents be
$m_1,\ldots,m_r$.  A rough degree bound for the unnormalized tower is
\[
 d=\prod_{i=1}^{r}m_i.
\]
A normal closure and every automorphism group can be bounded by a recurrence
using products, powers, and factorials.  For example, after replacing the
root-of-unity degree by an explicit bound $B_0$, one may use
\[
 B_{i+1}=B_i\,m_i^{B_i}
\]
and finally $B_r!$ as a permutation-table bound.  These estimates are
intentionally extravagant but primitive recursive.  They bound every
enumeration of embeddings, roots, automorphisms, subgroup closures, and
derived-series stages.  Efficiency is irrelevant to the foundational claim.

\section{The generic quintic has Galois group \texorpdfstring{$S_5$}{S5}}
\label{sec:generic-quintic}

The usual proof is often presented through the theorem on symmetric rational
functions.  That theorem is not needed here.  We define the coefficient field
directly as the subfield generated by the elementary symmetric functions and
identify the automorphism group by its action on the five roots.

\subsection{The fields and the polynomial}

Let $x_1,\ldots,x_5$ be algebraically independent indeterminates over $\Q$ and
put
\[
 L=\Q(x_1,x_2,x_3,x_4,x_5).
\]
For $1\le j\le5$, let
\[
 e_j=\sum_{1\le i_1<\cdots<i_j\le5}
       x_{i_1}\cdots x_{i_j}
\]
be the elementary symmetric functions, and let
\[
 K=\Q(e_1,e_2,e_3,e_4,e_5)\subseteq L.
\]
Define
\[
 P(T)=\prod_{i=1}^{5}(T-x_i).
 \tag{G1}\label{eq:generic-product}
\]
Expanding by elementary symmetric functions gives
\[
 P(T)=T^5-e_1T^4+e_2T^3-e_3T^2+e_4T-e_5\in K[T].
 \tag{G2}\label{eq:generic-coefficients}
\]

The elements $x_i$ are pairwise distinct.  Indeed, $x_i-x_j$ is a nonzero
polynomial when $i\ne j$, hence a nonzero element of the fraction field $L$.
Therefore $P$ is separable.  Each $x_i$ is algebraic over $K$ because it
satisfies the monic polynomial $P$, so the extension
$K(x_1,\ldots,x_5)/K$ is finite.  The polynomial splits in $L$, and
\[
 L=K(x_1,\ldots,x_5),
\]
so $L$ is the finite splitting field of $P$ over $K$.  Hence $L/K$ is normal
and separable, and its $K$-automorphism group is its Galois group.

All these claims have direct sparse-polynomial certificates.  Algebraic
independence is built into the representation of $L$ as the fraction field of
the polynomial ring $\Q[x_1,\ldots,x_5]$; no existence theorem for a generic
point is required.

\subsection{Every permutation gives an automorphism}

For $\pi\in S_5$, define a $\Q$-automorphism $\widehat\pi$ of $L$ by
\[
 \widehat\pi(x_i)=x_{\pi(i)}.
\]
On a rational function, apply the same variable renaming to its numerator and
denominator.  This is a primitive-recursive operation on sparse codes.
Because every $e_j$ is symmetric,
\[
 \widehat\pi(e_j)=e_j,
\]
so $\widehat\pi$ fixes $K$ pointwise.  Hence
\[
 S_5\longrightarrow\Gal(L/K),
 \qquad
 \pi\longmapsto\widehat\pi
 \tag{G3}\label{eq:S5-to-Gal}
\]
is a homomorphism.  It is injective: if $\pi(i)\ne\rho(i)$, then
$x_{\pi(i)}\ne x_{\rho(i)}$, so the two automorphisms differ on $x_i$.

\subsection{Every automorphism is a permutation of the roots}

Let $\sigma\in\Gal(L/K)$.  Since $\sigma$ fixes the coefficients of $P$,
\[
 P(\sigma(x_i))=\sigma(P(x_i))=0.
 \tag{G4}\label{eq:auto-root}
\]
But in $L$ we have the exact factorization~\eqref{eq:generic-product}.  Hence
\[
 \prod_{j=1}^{5}(\sigma(x_i)-x_j)=0.
\]
A field has no zero divisors, so $\sigma(x_i)=x_j$ for some $j$.  The root is
unique because the $x_j$ are pairwise distinct.  Thus there is a function
$\pi_\sigma:\{1,\ldots,5\}\to\{1,\ldots,5\}$ with
\[
 \sigma(x_i)=x_{\pi_\sigma(i)}.
\]
The function is injective because $\sigma$ is injective; a bounded finite
count proves that an injection on five elements is a bijection.  Therefore
$\pi_\sigma\in S_5$.

The automorphisms $\sigma$ and $\widehat{\pi_\sigma}$ agree on $\Q$ and on all
five generators $x_i$ of $L$, so they are equal.  Consequently
\eqref{eq:S5-to-Gal} is surjective as well as injective.

\begin{theorem}[Generic Galois group]
\label{thm:generic-galois-S5}
For the polynomial $P$ in~\eqref{eq:generic-coefficients},
\[
 \Gal(L/K)\cong S_5.
\]
The isomorphism and its inverse are given explicitly by variable renaming and
by reading the images of the five generators.
\end{theorem}

\begin{proof}
The preceding two subsections construct mutually inverse homomorphisms.  Every
step is an equality of rational functions or a bounded calculation on five
indices.
\end{proof}

\subsection{Why no fixed-field theorem is being hidden}

One may instead define $K$ as the fixed field $L^{S_5}$ and use Artin's
fixed-field theorem.  ProveIt's existing \code{GenericAbelRuffini.lean} follows
that elegant route.  For the iPRA arithmetization, the explicit field
$\Q(e_1,\ldots,e_5)$ is convenient because Theorem
\ref{thm:generic-galois-S5} needs only:

\begin{itemize}
\item sparse polynomial expansion of a five-factor product;
\item primitive-recursive renaming of five variables;
\item exact rational-function equality;
\item the no-zero-divisors rule; and
\item finite injection-to-bijection on a five-element list.
\end{itemize}

The fundamental theorem of symmetric polynomials, linear independence of
field embeddings, and degree counting are all unnecessary for this
identification.

\subsection{The same construction in every degree}

For any $n\ge2$, the identical argument with
$L_n=\Q(x_1,\ldots,x_n)$ and
$K_n=\Q(e_1,\ldots,e_n)$ gives
\[
 \Gal\left(L_n/K_n\right)\cong S_n
\]
for the generic degree-$n$ product.  The proof below only needs $n=5$ because
$S_5$ is already nonsolvable.  Higher degrees can therefore be handled either
by the generic $S_n$ construction or, more economically, by padding the
quintic.

\section{A fully finite proof that \texorpdfstring{$S_5$}{S5} is not solvable}
\label{sec:s5-nonsolvable}

A common proof invokes the simplicity of $A_5$.  For the present purpose that
is more than necessary.  We prove directly that $A_5$ is perfect, using one
explicit commutator.  This reduces the entire obstruction to calculations in
a group of $120$ permutations.

We compose permutations right-to-left and use
$[g,h]=ghg^{-1}h^{-1}$.

\subsection{The commutator subgroup of \texorpdfstring{$S_5$}{S5}}

The sign map
\[
 \sgn:S_5\longrightarrow\{\pm1\}
\]
has abelian target.  Every commutator lies in its kernel, so
\[
 [S_5,S_5]\subseteq A_5.
 \tag{S1}\label{eq:S5-derived-subset}
\]
On the other hand, a direct calculation gives
\[
 [(23),(12)]=(123).
 \tag{S2}\label{eq:S5-commutator}
\]
The commutator subgroup is normal.  Conjugating~\eqref{eq:S5-commutator} by
permutations shows that every $3$-cycle belongs to $[S_5,S_5]$.
Every even permutation is a product of $3$-cycles: write it as an even number
of transpositions and group those transpositions in pairs.  If two
transpositions share a letter, their product is a $3$-cycle; if they are
disjoint, for example,
\[
 (12)(34)=(123)(234).
\]
Hence $3$-cycles generate $A_5$, and
\[
 [S_5,S_5]=A_5.
 \tag{S3}\label{eq:S5-derived}
\]

\subsection{The commutator subgroup of \texorpdfstring{$A_5$}{A5}}

Both
\[
 a=(23)(45),
 \qquad
 b=(12)(45)
\]
are even.  A second direct calculation gives
\[
 [a,b]=(123).
 \tag{S4}\label{eq:A5-commutator}
\]
Thus $(123)\in[A_5,A_5]$.

Every $3$-cycle is conjugate to $(123)$ by an element of $A_5$.  To see this
constructively, first choose any permutation carrying the ordered triple
$(1,2,3)$ to the ordered triple defining the desired $3$-cycle.  If that
permutation is odd, compose it on the right with $(45)$.  This changes its
parity but leaves the images of $1,2,3$ unchanged.  Therefore an even
conjugator always exists.

The subgroup $[A_5,A_5]$ is normal in $A_5$, so it contains every $3$-cycle.
Since those cycles generate $A_5$,
\[
 [A_5,A_5]=A_5.
 \tag{S5}\label{eq:A5-perfect}
\]

\begin{theorem}[Nonsolvability of $S_5$]
\label{thm:S5-not-solvable}
The symmetric group $S_5$ is not solvable.
\end{theorem}

\begin{proof}
By~\eqref{eq:S5-derived}, the first derived subgroup is $A_5$.  By
\eqref{eq:A5-perfect}, every later derived subgroup is again $A_5$:
\[
 S_5^{(1)}=A_5,
 \qquad
 S_5^{(r)}=A_5\quad(r\ge1).
\]
Since $A_5$ has $60$ elements, the series never reaches the identity group.
\end{proof}

\subsection{Primitive-recursive verification}

There are only $5!=120$ permutations.  Encode a permutation by a five-entry
list, enumerate all $120$ lists that pass the bijection test, and compute:

\begin{enumerate}
\item the $60$ even permutations;
\item all commutators in $S_5$ and their generated subgroup;
\item all commutators in $A_5$ and their generated subgroup;
\item the two explicit identities~\eqref{eq:S5-commutator} and
  \eqref{eq:A5-commutator}.
\end{enumerate}

The supplied certificate script finds
\[
 |[S_5,S_5]|=60,
 \qquad
 [S_5,S_5]=A_5,
 \qquad
 |[A_5,A_5]|=60,
 \qquad
 [A_5,A_5]=A_5.
\]
These calculations are not needed as oracles: they can be reproduced by a
small primitive-recursive function and proved correct by induction on the
bounded enumeration.  The explicit human proof above makes the generated
subgroup certificates particularly short.

\section{The Abel--Ruffini contradiction}
\label{sec:abel-proof}

We now combine the three ingredients:

\begin{enumerate}
\item a polynomial split by a radical tower has solvable Galois group
  (Corollary~\ref{cor:split-rad-solvable});
\item the generic quintic has Galois group $S_5$
  (Theorem~\ref{thm:generic-galois-S5});
\item $S_5$ is not solvable
  (Theorem~\ref{thm:S5-not-solvable}).
\end{enumerate}

\subsection{No radical tower splits the generic quintic}

\begin{theorem}[Generic quintic obstruction]
\label{thm:generic-quintic-obstruction}
Let $K,L,P$ be as in Section~\ref{sec:generic-quintic}.  No radical tower over
$K$ contains all five roots $x_1,\ldots,x_5$ of $P$.
\end{theorem}

\begin{proof}
Suppose a radical tower contained all roots.  Since $L=K(x_1,\ldots,x_5)$,
the tower would contain the splitting field $L$.  By
Corollary~\ref{cor:split-rad-solvable}, $\Gal(L/K)$ would be solvable.  By
Theorem~\ref{thm:generic-galois-S5}, this group is isomorphic to $S_5$, which
is not solvable by Theorem~\ref{thm:S5-not-solvable}.  Contradiction.
\end{proof}

The proof is intuitionistically valid.  The assumption of a radical tower is
used to construct an explicit finite solvability certificate; the $S_5$
calculation refutes that certificate.  No double-negation elimination is used.

\subsection{From a universal formula to a radical tower}

Suppose there were a universal formula by radicals for monic quintics
\[
 T^5+a_1T^4+a_2T^3+a_3T^2+a_4T+a_5.
\]
The formula may contain finitely many branches, field operations, and radical
instructions.  Specialize its coefficient inputs to
\[
 a_1=-e_1,
 \quad a_2=e_2,
 \quad a_3=-e_3,
 \quad a_4=e_4,
 \quad a_5=-e_5
\]
in the field $K$.  Follow the branch guards, and adjoin the finitely many
radical witnesses in program order.  The resulting run is a radical tower.
Completeness of the formula supplies all five roots and therefore a complete
factorization of $P$.  This contradicts
Theorem~\ref{thm:generic-quintic-obstruction}.

This specialization argument covers formulas with exceptional cases.  A
formula that is truly universal must have a valid branch for the generic
coefficient tuple.  It also covers formulas whose output roots are unordered,
whose branches repeat roots, or whose radicals have several possible values:
any successful run still generates a radical tower containing a complete root
list.

\begin{theorem}[Abel--Ruffini in degree five]
There is no universal complete formula, built from rational constants, field
operations, finite equality case distinctions, and extraction of finitely
many radicals, that solves every degree-five equation over characteristic-zero
fields.
\end{theorem}

\subsection{Every degree greater than five}

For $n\ge5$, put
\[
 P_n(T)=T^{n-5}P(T).
\]
This is monic of exact degree $n$.  Its root multiset contains
$x_1,\ldots,x_5$ together with $n-5$ zero roots.  If a radical tower split
$P_n$, it would contain all $x_i$ and hence split $P$, contradicting
Theorem~\ref{thm:generic-quintic-obstruction}.

\begin{corollary}[Abel--Ruffini in every degree at least five]
For each $n\ge5$, no universal complete radical formula exists for all
polynomials of degree $n$.
\end{corollary}

The padding argument deliberately proves a statement about complete solution.
The padded polynomial has the easy root zero; producing that one root does not
solve the equation.

\subsection{Why the coefficient field is legitimate}

A reader may ask whether a formula intended for complex coefficients must be
tested over the rational-function field $K$.  There are two equivalent
answers.

\paragraph{Formal-expression answer.}
A universal radical formula is an identity in the coefficient indeterminates.
Its operations are algebraic and therefore make sense over every
characteristic-zero field where its denominators are nonzero and its guarded
branches are followed.  The generic field is precisely the universal test
object for such identities.

\paragraph{Specialization answer.}
If one insists on complex coefficients, embed a finitely generated
specialization of the generic data into a sufficiently large characteristic-
zero field and preserve the finitely many nonzero denominators.  A successful
universal formula would then yield the same radical-tower consequence.  An
alternative route, already present in ProveIt, uses explicit rational
polynomials with nonsolvable Galois group.  Section~\ref{sec:explicit-quintic}
records one elementary certificate route.

\subsection{Strong and weak formulations}

The following implications are useful when comparing formal statements.

\[
\begin{array}{c}
\text{a fixed coefficient formula returns all roots}\\[0.2em]
\Downarrow\\[-0.1em]
\text{each input polynomial has a complete radical expression}\\[0.2em]
\Downarrow\\[-0.1em]
\text{the generic quintic has a complete radical factorization}\\[0.2em]
\Downarrow\\[-0.1em]
\Gal(P/K)\text{ is solvable}.
\end{array}
\]

The proof refutes the bottom statement.  Therefore it refutes every stronger
notion above it, including nonuniform per-root expressions when they are
required for all five roots.  This is why the argument is insensitive to the
surface syntax chosen for ``formula by radicals.''

\section{Internalizing the proof in intuitionistic PRA}
\label{sec:internalization}

This section turns the preceding mathematics into a precise arithmetization
plan.  The goal is not merely to say that the proof is ``finitary.''  We
identify the primitive-recursive functions, the certificate relations, the
induction measures, and the final iPRA sentences.

\subsection{The base syntax}

Fix natural-number codes for the following datatypes.

\begin{align*}
 \code{ZCode} & &&\text{signed integers},\\
 \code{QCode} & &&\text{rational numbers},\\
 \code{MonomCode}(v) & &&\text{monomials in $v$ variables},\\
 \code{PolyCode}(v) & &&\text{normalized sparse polynomials},\\
 \code{RatFunCode}(v) & &&\text{fractions of polynomial codes},\\
 \code{PermCode}(d) & &&\text{permutations of $d$ symbols},\\
 \code{GroupCode} & &&\text{finite group tables},\\
 \code{ExprCode} & &&\text{field and radical expressions},\\
 \code{ProofCode} & &&\text{finite derivations},\\
 \code{TowerCode} & &&\text{radical-tower certificates}.
\end{align*}

The constructor tags, list projections, and all recursors are primitive
recursive.  Parsing is total: malformed codes are interpreted as a fixed
error object, so every checker is a total natural-number function.

\subsection{Polynomial and rational-function normalization}

Define the following primitive-recursive functions.

\begin{align*}
 \operatorname{pnorm}&:\code{PolyCode}(v)\to\code{PolyCode}(v),\\
 \operatorname{padd},\operatorname{pmul},\operatorname{pneg}
  &:\code{PolyCode}(v)^2\to\code{PolyCode}(v),\\
 \operatorname{psubst}&:\code{PolyCode}(v)\times
   \code{PolyCode}(w)^v\to\code{PolyCode}(w),\\
 \operatorname{rfeq}&:\code{RatFunCode}(v)^2\to\{0,1\}.
\end{align*}

The sparse normalizer sorts monomials by a primitive-recursive lexicographic
rank, combines adjacent equal exponent vectors, and deletes zero coefficients.
Correctness is proved by quantifier-free induction on the source list.  For
fractions $(P,Q)$ and $(R,S)$,
\[
 \operatorname{rfeq}((P,Q),(R,S))=1
 \quad\Longleftrightarrow\quad
 \operatorname{pnorm}(PS-RQ)=0,
\]
subject to checked nonzero denominators.  This gives decidable equality in the
generic rational-function field $L$.

The elementary symmetric functions $e_1,\ldots,e_5$ and the product
$\prod_i(T-x_i)$ are fixed sparse codes.  A one-time normalization certificate
proves equation~\eqref{eq:generic-coefficients}.

\subsection{Low-degree formula checkers}

For each degree $n\le4$, define a total function
\[
 \operatorname{lowCheck}_n(a,w)\in\{0,1\},
\]
where $a$ is a coefficient tuple and $w$ contains the nonzero and radical
witnesses.  It checks:

\begin{itemize}
\item the leading coefficient is nonzero;
\item every radical power equation;
\item every denominator side condition;
\item the coefficient-by-coefficient factorization identity.
\end{itemize}

The root-list constructor
\[
 \operatorname{lowRoots}_n(a,w)
\]
is primitive recursive.  The core iPRA theorems are equations of the form
\[
 \operatorname{lowCheck}_n(a,w)=1
 \to
 \operatorname{factorCheck}
   \bigl(a,\operatorname{lowRoots}_n(a,w)\bigr)=1.
 \tag{I1}\label{eq:ipra-low}
\]
For $n=2$ the proof is a single normalized identity.  For $n=3$ it uses the
three relations in Section~\ref{sec:cubic}.  For $n=4$ it composes the Ferrari
and quadratic identities.  No induction is needed beyond correctness of the
normalizer.

A separate primitive-recursive branch constructor implements the compatible
radical choices:

\begin{itemize}
\item in the cubic, use the supplied coefficient guard $p=0$ or $p\ne0$;
  the $p=0$ branch adjoins two cube roots and deduces $uv=0$ from
  $(uv)^3=0$;
\item in the quartic, use the supplied coefficient guard $q=0$ or $q\ne0$;
  the first branch is biquadratic and the second proves the Ferrari radical
  $s$ nonzero before division;
\item carry the selected equality or nonzero derivation into the denominator
  checker.
\end{itemize}

Because the guards are values of \code{rfeq}, these are intuitionistically
legitimate case distinctions.

\subsection{Finite permutation groups}

Define
\[
 \operatorname{allPerm}(d)
\]
by bounded enumeration of the $d^d$ lists of length $d$ and filtering by the
bijection predicate.  This is inefficient but primitive recursive.  For
$d=5$ the output has length $120$.  Define composition, inverse, sign,
commutator, subgroup closure, and derived subgroup by bounded folds.

Let \code{S5} and \code{A5} be the resulting fixed codes.  The following are
closed equations between numerals and primitive-recursive terms:
\begin{align}
 |\code{S5}|&=120,\\
 |\code{A5}|&=60,\\
 \operatorname{derived}(\code{S5})&=\code{A5},\\
 \operatorname{derived}(\code{A5})&=\code{A5}.
 \tag{I2}\label{eq:ipra-s5}
\end{align}
They can be proved either by evaluation of the full tables or by the much
smaller explicit commutator certificates of Section~\ref{sec:s5-nonsolvable}.
The latter route is preferable for proof size: verify two products, normality
of a generated subgroup, conjugacy of $3$-cycles, and the generation of $A_5$
by $3$-cycles.

Define the Boolean predicate
\[
 \operatorname{solvableB}(G)=1
\]
by iterating the derived-subgroup function $|G|$ times and testing whether the
identity subgroup has been reached.  Equation~\eqref{eq:ipra-s5} proves
\[
 \operatorname{solvableB}(\code{S5})=0.
 \tag{I3}\label{eq:S5-bool-false}
\]

\subsection{Coding radical-tower solvability certificates}

A tower certificate contains:

\begin{enumerate}
\item a finite list of exponents and radical power relations;
\item the finite lists of conjugate radicands used in each normal-closure
  stage;
\item root-of-unity data;
\item a finite automorphism table for each normal stage;
\item restriction maps and exactness witnesses;
\item embeddings of each abelian kernel in a product of cyclic groups;
\item a final restriction map onto the claimed splitting-field group.
\end{enumerate}

The local checker \code{towerCheck} verifies every polynomial relation and
every finite table equation.  The function
\[
 \operatorname{seriesOfTower}(t)
\]
reads the exact sequences and concatenates the derived-series certificates of
the abelian kernels and quotients.  Its correctness theorem has the form
\[
 \operatorname{towerCheck}(t)=1
 \to
 \operatorname{solvableSeriesCheck}
    (\operatorname{seriesOfTower}(t))=1.
 \tag{I4}\label{eq:tower-to-series}
\]
The proof uses quantifier-free induction on the tower length; the invariant is
an equation saying that the certificate generated for the prefix passes its
checker.  No semantic induction over arbitrary fields occurs inside iPRA.

The extension-of-embeddings steps are also proof-carrying.  A certificate lists
the image chosen for each new generator among a bounded finite root list.  The
checker verifies its defining polynomial and every relation to earlier
generators.  Thus no unbounded search or choice operation is trusted.

\subsection{Coding the generic splitting field action}

The arithmetized generic proof needs only rational-function syntax and the
fixed set of $120$ permutations.

For each permutation code $\pi$, define
\[
 \operatorname{rename}_\pi:\code{RatFunCode}(5)\to\code{RatFunCode}(5)
\]
by renaming exponent-vector coordinates.  The fixed polynomial codes for
$e_j$ satisfy
\[
 \operatorname{rfeq}(\operatorname{rename}_\pi(e_j),e_j)=1
\]
for all $\pi$ and $j$, by a bounded check.  Hence every permutation fixes the
coefficient presentation.

Conversely, an automorphism certificate of the splitting presentation lists
the images of $x_1,\ldots,x_5$.  The factorization code for $P$ and the field
cancellation rules prove that each image is one of the five variable codes.
Distinctness is checked by normalizing $x_i-x_j$.  A bounded injection test
therefore recovers a unique permutation.  The two transformations are inverse
by five rational-function equalities.

The resulting theorem is an isomorphism certificate
\[
 \operatorname{galCode}(P)=\code{S5}.
 \tag{I5}\label{eq:gal-code-s5}
\]
No fixed-field quantification is required in this coding.

\subsection{The arithmetized contradiction}

Let \code{generic5} be the fixed code of $P$, and let
\code{RadSplit(generic5,t)} be the formula asserting that $t$ is a checked
radical splitting certificate.  The radical-tower compiler and
\eqref{eq:tower-to-series} give
\[
 \operatorname{RadSplit}(\code{generic5},t)
 \to
 \operatorname{solvableB}(\operatorname{galCode}(P))=1.
\]
Using~\eqref{eq:gal-code-s5}, this implies
$\operatorname{solvableB}(\code{S5})=1$, contradicting
\eqref{eq:S5-bool-false}.  Therefore iPRA proves
\[
 \forall t\;\neg\operatorname{RadSplit}(\code{generic5},t).
 \tag{I6}\label{eq:ipra-no-radsplit}
\]

A primitive-recursive specialization compiler \code{compileGeneric} maps a
purported universal degree-five formula code $c$ and a successful run witness
$w$ to a radical splitting certificate for \code{generic5}.  Hence
\[
 \forall c\;\bigl(
   \operatorname{UniversalRadicalFormula}(5,c)
   \to\bot\bigr).
 \tag{I7}\label{eq:ipra-abel5}
\]
Padding is a primitive-recursive polynomial operation, so iPRA proves
\[
 \forall n\;\bigl(5\le n\to
   \forall c\;\neg\operatorname{UniversalRadicalFormula}(n,c)\bigr).
 \tag{I8}\label{eq:ipra-abel-all}
\]

\subsection{Transfer to the concise native calculus}

The final formulas~\eqref{eq:ipra-low}, \eqref{eq:ipra-abel5}, and
\eqref{eq:ipra-abel-all} are ordinary iPRA sentences.  The ProveIt file
\code{NativeBiInterpretation.lean}
defines translations between ordinary iPRA formulas and native formulas and
proves \code{native_iff_pra}; the round trips
are provably the identity up to the established equivalence.  Therefore a
Lean proof of the ordinary iPRA derivation immediately yields a derivation in
the selected concise calculus.  No new foundational argument about the tiny
presentation is required.

\subsection{No hidden use of stronger arithmetic}

The proof uses large finite bounds, but proof-theoretic strength is not the
same as computational efficiency.  Factorial, bounded exponentiation, sparse
normalization, finite Gaussian elimination, and finite subgroup closure are
all primitive recursive.  Their correctness proofs use induction on explicit
counters with decidable invariants.  The proof does not require:

\begin{itemize}
\item induction for arbitrary formulas;
\item unbounded minimization;
\item transfinite induction;
\item compactness, Zorn's lemma, or a global algebraic closure;
\item excluded middle for arbitrary field-theoretic propositions.
\end{itemize}

The use of splitting fields and normal closures in the human proof is shorthand
for finite witness records.  In the arithmetized proof, the witnesses are
present in the certificate and every property used later is locally checked.

\section{Lean formalization plan}
\label{sec:lean-plan}

\subsection{Formalization status}

There are two distinct layers.

\begin{enumerate}
\item The \emph{semantic algebra layer} is already kernel-checked in the
  inspected ProveIt tree.  It contains the field-generic factorization
  theorems for degrees at most four, the typed radical-expression language,
  the semantic radical closure, explicit rational Abel--Ruffini examples, and
  the generic fixed-field construction.
\item The \emph{iPRA arithmetization layer} specified in
  Sections~\ref{sec:ipra-framework} and~\ref{sec:internalization} is not yet a
  completed Lean development.  The present archive gives its definitions,
  theorem interfaces, proof decomposition, and an exact finite checker, but
  it does not claim an end-to-end kernel-checked theorem of type
  \code{PRAAx derives AbelRuffini}.
\end{enumerate}

This distinction matters.  Importing a classical field-theory theorem into
Lean proves the mathematics, but it does not by itself prove that a derivation
exists in the weak object theory iPRA.  The missing bridge must internalize the
finite codes and the correctness proofs of their primitive-recursive
checkers.

\subsection{Existing files to reuse}

The following imports should remain the semantic reference implementation.

\begin{longtable}{p{0.34\textwidth}p{0.58\textwidth}}
\toprule
Existing module & Reused theorem families\\
\midrule
\endhead
\code{PolynomialFormulas.Basic} &
\code{linear_eq_zero_iff}, \code{quadratic_formula_factorization},
\code{quadratic_eq_zero_iff}.\\
\code{PolynomialFormulas.Cubic} &
\code{depress_monic_cubic}, \code{cardano_depressed},
\code{cardano_factorization}, \code{cubic_factorization},
\code{cubic_eq_zero_iff}.\\
\code{PolynomialFormulas.Quartic} &
\code{depress_monic_quartic}, \code{ferrari_factorization},
\code{solveDepressedQuartic_factorization},
\code{quartic_factorization}, \code{quartic_eq_zero_iff}.\\
\code{PolynomialFormulas.AbelRuffini} &
\code{RadicalExpression},
\code{nonempty_iff_mem_solvableByRad}, and the explicit rational quintic
obstruction.\\
\code{PolynomialFormulas.GenericAbelRuffini} &
The rational-function-field model, root orbit, splitting-field construction,
Galois equivalence with a symmetric group, and rootwise obstruction.\\
\code{ConciseIPRA.NativeBiInterpretation} &
\code{native_iff_pra}, \code{pra_iff_native}, and the two round-trip
translation theorems.\\
\bottomrule
\end{longtable}

The file \code{PolynomialFormulas.Audit} already asks Lean to print the axioms
of the main semantic theorems.  The new development should add a separate
audit for the iPRA endpoints and fail continuous integration if any endpoint
contains \code{sorryAx} or an unintended axiom.

\subsection{Proposed module tree}

A practical directory is
\[
 \code{Logic/PrimitiveRecursiveArithmetic/Applications/AbelRuffini/}.
\]
The following modules give a dependency order with no cycles.

\begin{description}[style=nextline]
\item[\code{Codes.lean}]
Pairing, bounded lists, zig-zag integers, rationals, exponent vectors, sparse
polynomials, rational functions, permutations, finite tables, and generic
proof-tree codes.  Each operation is accompanied by a \code{Primrec} theorem.

\item[\code{SparsePolynomial.lean}]
Canonical sorting and coefficient collection, evaluation semantics, and the
fundamental theorem
\[
 \operatorname{eval}(\operatorname{pnorm}(p),\rho)
 =\operatorname{eval}(p,\rho).
\]
Prove the list-loop invariants first in ordinary Lean; then reflect them into
\code{PRAAx} using the existing primitive-recursion infrastructure.

\item[\code{RationalFunction.lean}]
Nonzero-denominator records, cross-multiplication equality, variable renaming,
and substitution.  This module proves the exact identities for the five
symmetric generators and the generic product polynomial.

\item[\code{LowDegreeCertificates.lean}]
Codes for the quadratic, Cardano, biquadratic, and Ferrari witness packages;
root-list constructors; and Boolean factorization checkers.  The semantic
correctness lemmas can be connected to the existing factorization theorems,
while the object-theory results use only the sparse-normalizer equations.

\item[\code{FinitePermutationGroups.lean}]
Bounded enumeration of \code{Equiv.Perm (Fin d)}, table construction,
commutators, subgroup closure, derived series, and
\code{solvableB}.  Include compact certificates for
$[S_5,S_5]=A_5$ and $[A_5,A_5]=A_5$ rather than relying only on a large
\code{native_decide} computation.

\item[\code{RadicalTowerCertificate.lean}]
The proof-carrying finite-extension and radical-layer records, local checker,
normal-layer compiler, exact-sequence certificates, and the theorem that every
accepted tower produces an accepted solvable-series certificate.

\item[\code{GenericQuinticCertificate.lean}]
Codes for $e_1,\ldots,e_5$ and $P$; permutation renaming; the map from a
permutation to an automorphism table; the inverse map reading the images of
$x_i$; and the checked isomorphism with \code{S5Code}.

\item[\code{AbelRuffiniPRA.lean}]
Definitions of \code{RadSplit}, \code{UniversalRadicalFormula}, and the
generic specialization compiler.  Combine the tower-solvability theorem,
generic Galois certificate, and $S_5$ computation into the ordinary iPRA
derivations.

\item[\code{AbelRuffiniNative.lean}]
Translate the ordinary formulas to the selected native calculus and discharge
them with \code{native_iff_pra}.  This file should contain no algebra.

\item[\code{Audit.lean}]
\texttt{\#check} and \texttt{\#print axioms} for every public endpoint, plus closed
examples for the degree-two, degree-three, degree-four, and $S_5$ checkers.
\end{description}

\subsection{Suggested Lean interfaces}

The exact first-order syntax names should follow ConciseIPRA, but the public
mathematical interfaces should look like the following.  These declarations
are schematic: the archive's \code{Lean/IPRAAbelRuffiniBlueprint.lean} records
the module boundaries without pretending that the missing proofs have already
been checked.

\begin{lstlisting}[style=lean,language={},caption={Certificate-level endpoints}]
structure LowDegreeWitness (n : Fin 5) where
  coeffs : List QCode
  data   : ExprCode
  proof  : ProofCode

def lowCheck (n : Fin 5) : LowDegreeWitness n -> Bool
def lowRoots (n : Fin 5) : LowDegreeWitness n -> List ExprCode

theorem lowCheck_sound {n} {w : LowDegreeWitness n} :
    lowCheck n w = true -> factorCheck w.coeffs (lowRoots n w) = true

def towerCheck : TowerCode -> Bool
def solvableB : GroupCode -> Bool

theorem towerCheck_implies_solvable {t} :
    towerCheck t = true -> solvableB (splittingGroup t) = true

theorem genericQuintic_galCode :
    genericQuinticGroup = S5Code

theorem S5_not_solvableB : solvableB S5Code = false

theorem no_generic_radical_split :
    forall t, RadicalSplit genericQuintic t -> False
\end{lstlisting}

The object-theory endpoint then has the shape

\begin{lstlisting}[style=lean,language={},caption={Object-theory endpoints, schematic}]
theorem pra_lowDegree_factorization (n : Fin 5) :
  [] |-PRA- lowDegreeSentence n

theorem pra_abelRuffini_degree_five :
  [] |-PRA- not (existsCode (UniversalRadicalFormula 5))

theorem pra_abelRuffini_degree_ge_five :
  [] |-PRA- forallNat fun n =>
    5 <= n -> not (existsCode (UniversalRadicalFormula n))

theorem native_abelRuffini_degree_ge_five :
  NPrf [] (NFm.ofPRA abelRuffiniSentence)
\end{lstlisting}

Here \code{|-PRA-} abbreviates ProveIt's actual notation
\(\Gamma \vdash[\mathrm{PRAAx}] A\).  The final proof is a one-line use of
\code{native_iff_pra} after the ordinary derivation exists.

\subsection{How to internalize a primitive-recursive checker}

For each checker, use the same four-stage pattern.

\begin{enumerate}
\item Define a total Lean function on natural-number codes.  Malformed input
  returns \code{false} or a fixed error code.
\item Prove the function primitive recursive compositionally.  Do not mark it
  recursive by an opaque kernel shortcut.
\item State its loop invariant as an equality between primitive-recursive
  functions.  Prove that invariant in \code{PRAAx} by quantifier-free
  induction, using the equations for primitive recursion already available in
  ConciseIPRA.
\item Derive the externally readable semantic theorem from the accepted
  certificate.  The semantic theorem may use structures and fields; the
  object-theory theorem mentions only the code and checker equations.
\end{enumerate}

For fixed identities such as the Cardano and Ferrari factorizations, an even
smaller route is available: generate a closed proof trace of sparse-ring
normalization and verify that trace with a primitive-recursive step checker.
The iPRA proof then establishes soundness of the generic step checker once and
applies it to a concrete numeral coding the trace.

\subsection{Milestones and completion criterion}

A clear path to completion is:

\begin{enumerate}
\item \textbf{Low-degree milestone.}  Kernel-check
  \code{pra_lowDegree_factorization} for $n=1,2,3,4$.  This validates the
  coding and normalization layer before any Galois theory is internalized.
\item \textbf{Finite-group milestone.}  Kernel-check
  \code{S5_not_solvableB} with no native oracle in the public proof term.
\item \textbf{Generic-action milestone.}  Kernel-check the mutually inverse
  maps between the $120$ renamings and the automorphism certificate.
\item \textbf{Radical-tower milestone.}  Kernel-check the compiler from an
  accepted radical tower to a solvable-series certificate.
\item \textbf{Ordinary iPRA endpoint.}  Kernel-check the degree-five
  contradiction and the padding theorem in \code{PRAAx}.
\item \textbf{Native endpoint.}  Transfer through
  \code{native_iff_pra}, run \texttt{\#print axioms}, and place the result under
  the repository's existing no-placeholder audit.
\end{enumerate}

The formalization is complete only when the last two endpoints print no
\code{sorryAx}, no extra arithmetic axiom, and no unrecorded computation
oracle.  The Python/SymPy script in this archive is useful regression evidence
but is not part of that trust claim.

\section{Optional rational specialization: \texorpdfstring{$X^5-X-1$}{X5-X-1}}
\label{sec:explicit-quintic}

The generic proof is the cleanest route for the iPRA architecture because its
Galois action is literal variable permutation.  For readers who want a single
polynomial over $\Q$, this section gives a compact finite certificate for
\[
 f(X)=X^5-X-1.
 \tag{E1}\label{eq:selmer5}
\]
The only additional theorem needed is the finite Dedekind--Frobenius lemma:
for a prime not dividing the discriminant, the degrees of the distinct
irreducible factors modulo that prime occur as the cycle lengths of an element
of the rational Galois group.  That lemma can itself be arithmetized, but doing
so adds more local algebra than the generic proof requires.

\subsection{Irreducibility modulo three}

Modulo $3$, the values at $0,1,2$ are all $2$, so $f$ has no linear factor.
The monic irreducible quadratics over $\F_3$ are
\[
 X^2+1,
 \qquad X^2+X+2,
 \qquad X^2+2X+2.
\]
Polynomial division gives the respective remainders
\[
 -1,
 \qquad X-1,
 \qquad X-1.
\]
Thus none divides $f$.  A reducible polynomial of degree five has a factor of
degree one or two: among the degrees of its nonconstant factors, the smallest
is at most $2$.  Hence $f$ is irreducible over $\F_3$, and Gauss's lemma gives
irreducibility over $\Q$.

The same modular factorization is a single irreducible factor of degree five,
so the Dedekind--Frobenius lemma supplies a $5$-cycle in the Galois group.

\subsection{A transposition from reduction modulo two}

Modulo $2$ one has the exact factorization
\[
 X^5+X+1
 =(X^2+X+1)(X^3+X^2+1).
 \tag{E2}\label{eq:mod2-factor}
\]
The quadratic has no root in $\F_2$, and the cubic also has no root in
$\F_2$; hence both factors are irreducible.  A resultant computation gives
\[
 \operatorname{disc}(X^5-X-1)=2869=19\cdot151.
 \tag{E3}\label{eq:selmer-disc}
\]
In particular, neither $2$ nor $3$ divides the discriminant.  Reduction
modulo $2$ is therefore square-free, and the factor degrees $2$ and $3$ give
an element of cycle type $(2)(3)$ in the Galois group.  The cube of this
element is a transposition.

\subsection{The Galois group is \texorpdfstring{$S_5$}{S5}}

Let $c$ be the $5$-cycle obtained modulo $3$ and let $t$ be the transposition
obtained from the cube of the $(2)(3)$ element.  The conjugates
\[
 c^k t c^{-k},\qquad k=0,1,2,3,4,
\]
are transpositions whose edges form a connected graph on the five letters.
Indeed, if $t$ exchanges vertices at cyclic distance $d$, then the conjugates
join $i$ to $i+d$ modulo $5$; because $5$ is prime and $d\ne0$, this graph is
a $5$-cycle.  Transpositions along the edges of any connected graph generate
the full symmetric group.  Hence
\[
 \Gal(f/\Q)=S_5.
\]
By Theorem~\ref{thm:S5-not-solvable}, this group is not solvable.

\begin{corollary}[Explicit rational Abel--Ruffini witness]
No root of $X^5-X-1$ is expressible by radicals over $\Q$.
\end{corollary}

\begin{proof}
Suppose one root $\alpha$ lay in a radical tower.  Take a finite solvable
normal closure $N/\Q$ of that tower as in
Theorem~\ref{thm:radical-group-solvable}.  Since $f$ is irreducible, every
conjugate root is the image of $\alpha$ under a $\Q$-embedding; normality of
$N$ places all those images in $N$.  Thus the splitting field of $f$ is
contained in $N$, and its Galois group is a quotient of the solvable group
$\Gal(N/\Q)$.  This contradicts $\Gal(f/\Q)=S_5$.
\end{proof}

This rational example strengthens the coefficient-field interpretation of
Section~\ref{sec:abel-proof}.  It is not needed for the core generic proof, but
it is useful as an independently checkable specialization and aligns with the
rootwise rational results already developed in ProveIt.

\section{Executable certificate checks}
\label{sec:checks}

The archive contains
\[
 \code{certificates/verify_certificates.py}.
\]
It performs exact symbolic and finite checks; it uses no floating-point
arithmetic.  From the archive root, run

\begin{lstlisting}[style=lean,language={}]
python3 certificates/verify_certificates.py
\end{lstlisting}

The script writes \code{verification_report.json} and exits with a nonzero
status if a required check fails.  Its role is reproducibility and error
detection.  It is not substituted for the human proof and is not part of the
Lean kernel trust base.

\subsection{Checks performed}

\begin{longtable}{p{0.48\textwidth}p{0.18\textwidth}p{0.25\textwidth}}
\toprule
Check & Result & Mathematical role\\
\midrule
\endhead
Quadratic factorization reduced by
$s^2=b^2-4ac$ & $0$ & Equation~\eqref{eq:quad-factor}.\\
Cardano three-factor product reduced by
$\omega^2+\omega+1$, $uv=-p/3$, and $u^3+v^3=-q$ & $0$ &
Equation~\eqref{eq:cardano-factorization}.\\
Ferrari two-quadratic product reduced by the three parameter equations & $0$ &
Equation~\eqref{eq:ferrari-target}.\\
Denominator-free resolvent plus eight times the monic resolvent & $0$ &
Equivalence of~\eqref{eq:ferrari-resolvent} and
\eqref{eq:ferrari-resolvent-monic}.\\
Order of $S_5$ and $A_5$ & $120$, $60$ & Finite group coding.\\
Derived subgroup of $S_5$ & $A_5$ of order $60$ &
Equation~\eqref{eq:S5-derived}.\\
Derived subgroup of $A_5$ & $A_5$ of order $60$ &
Equation~\eqref{eq:A5-perfect}.\\
Two displayed commutators & both $(123)$ &
Equations~\eqref{eq:S5-commutator} and~\eqref{eq:A5-commutator}.\\
Discriminant of $X^5-X-1$ & $2869$ &
Equation~\eqref{eq:selmer-disc}.\\
Factorization modulo $2$ & irreducible degrees $2,3$ &
A $(2)(3)$ Frobenius cycle.\\
Factorization modulo $3$ & one irreducible degree-$5$ factor &
A $5$-cycle and rational irreducibility.\\
\bottomrule
\end{longtable}

The generated report records
\[
 \code{"all_required_checks_pass": true}.
\]

\subsection{Why the formal proof should not call SymPy}

The Groebner reductions in the script are convenient independent checks, but
they would be an unnecessarily large trust dependency in an iPRA proof.  The
formal route is smaller:

\begin{itemize}
\item use a verified sparse normalizer for commutative-ring expressions;
\item store a short reduction trace using the named radical relations;
\item verify each trace step with a primitive-recursive checker;
\item prove the generic checker sound once by quantifier-free induction.
\end{itemize}

Likewise, the full $120^2$ commutator enumeration is useful regression
evidence, whereas the proof term can use the two explicit commutators and the
short generation argument from Section~\ref{sec:s5-nonsolvable}.

\subsection{Reproducibility boundary}

The exact report in this archive was produced with the installed Python and
SymPy environment.  Future library versions may print factor lists in a
different order, but the script's required Boolean checks depend only on exact
zero remainders, subgroup equality, commutator equality, and the integer
discriminant.  The source and generated JSON are both included so that the
calculation can be inspected without trusting this prose summary.

\section{Further research questions}
\label{sec:further-research}

The proof architecture suggests several concrete follow-up projects.

\subsection{Finish the object-theory bridge}

The highest-priority task is to implement the module tree of
Section~\ref{sec:lean-plan} and produce the first end-to-end theorem
\[
 []\vdash_{\ipra}
 \forall c\;\neg\operatorname{UniversalRadicalFormula}(5,c).
\]
A useful engineering question is whether the ordinary \code{PRAAx} proof or a
direct proof in one of ConciseIPRA's smaller clause calculi yields the smaller
checked certificate.  The bi-interpretation guarantees equivalence, but proof
sizes may differ substantially.

\subsection{Replace automorphism tables by permutation presentations}

The radical-normal-closure certificate proposed here can list finite
multiplication tables.  Those tables may be enormous.  A more economical
certificate would present each group as a permutation group generated by a
short list and prove solvability through a subnormal series with abelian factor
certificates.  All membership tests remain bounded and primitive recursive.
The research problem is to find the smallest checker and certificate format
whose soundness still has a transparent iPRA proof.

\subsection{A branch-free constructive quartic object}

The report uses the decidable coefficient guard $q=0$ versus $q\ne0$.
One can instead build a universal finite splitting algebra in which the
Ferrari parameter equations are imposed simultaneously, then prove a product
identity without selecting a branch.  When is that algebra guaranteed to be a
nonzero finite product of fields constructively?  A satisfactory answer could
yield a formula theorem over fields without decidable equality while retaining
an explicit computational interpretation.

\subsection{A rational counterexample in iPRA}

Section~\ref{sec:explicit-quintic} reduces the rational example
$X^5-X-1$ to finite modular arithmetic plus the Dedekind--Frobenius cycle
lemma.  Formalizing that lemma in iPRA would give a closed integer-coefficient
counterexample, avoiding rational-function fields in the final sentence.
ProveIt's existing Selmer development supplies a strong semantic target.  A
particularly attractive endpoint is a numeral-only certificate proving that
no root of the displayed quintic has a radical-expression code.

\subsection{Proof-size lower and upper bounds}

The mathematical proof is elementary, but its naively expanded arithmetic
certificate may be very large.  Natural quantitative questions are:

\begin{itemize}
\item What is the smallest sparse-ring trace for the Cardano and Ferrari
  identities?
\item Is the explicit perfectness proof of $A_5$ smaller than evaluation of a
  generated subgroup table after compression?
\item How much does a subnormal-series certificate reduce the radical-tower
  proof compared with a derived-series certificate?
\item What is the compressed size of the final proof under the current
  ConciseIPRA checker metric?
\end{itemize}

These are mathematical proof-compression questions, not merely implementation
tuning: a better certificate format can expose a simpler invariant and lead to
a genuinely shorter iPRA derivation.

\subsection{Solvability decisions beyond degree five}

ProveIt already contains primitive-recursive or recursive decision procedures
for individual quintics and sextics.  The present generic obstruction and
those decision procedures could be connected at the code level: the same
\code{RadicalExpression} semantics should be reflected both by a positive
solver certificate and by a negative Galois-group certificate.  A unified
interface would make it possible to return one of two proof objects for an
input polynomial: a complete radical factorization or a checked nonsolvable
Galois-group witness.

\section{References and source map}
\label{sec:references}

\begin{thebibliography}{99}

\bibitem{Abel1824}
Niels Henrik Abel,
\emph{Mémoire sur les équations algébriques où l'on démontre l'impossibilité
de la résolution de l'équation générale du cinquième degré}, 1824.

\bibitem{Ruffini}
Paolo Ruffini,
\emph{Teoria generale delle equazioni}, 1799, and subsequent editions.

\bibitem{ArtinGalois}
Emil Artin,
\emph{Galois Theory}.
The standard finite-extension proof that radical extensions have solvable
Galois groups is reorganized in Section~\ref{sec:radical-solvable} into
proof-carrying finite layers.

\bibitem{DummitFoote}
David S. Dummit and Richard M. Foote,
\emph{Abstract Algebra}, chapters on solvable groups, field extensions, and
solvability by radicals.

\bibitem{LangAlgebra}
Serge Lang,
\emph{Algebra}, sections on Galois theory and equations by radicals.

\bibitem{HajekPudlak}
Petr Hájek and Pavel Pudlák,
\emph{Metamathematics of First-Order Arithmetic}.
This is a standard reference for coding syntax and primitive-recursive proof
predicates in weak arithmetic.

\bibitem{TroelstraVanDalen}
A. S. Troelstra and D. van Dalen,
\emph{Constructivism in Mathematics}.
The witness-based treatment of disjunction, existence, and partial algebraic
operations follows the constructive discipline emphasized in this report.

\bibitem{ProveItPolynomialFormulas}
Vladimir Reshetnikov,
\emph{ProveIt: Polynomial formulas and the Abel--Ruffini obstruction},
\url{https://github.com/VladimirReshetnikov/ProveIt/tree/25da87848ad4042c318af9bbc7f2e1c24614d6c3/Algebra/PolynomialFormulas}.
The relevant Lean files are \code{Basic.lean}, \code{Cubic.lean},
\code{Quartic.lean}, \code{AbelRuffini.lean}, and
\code{GenericAbelRuffini.lean}.

\bibitem{ProveItConciseIPRA}
Vladimir Reshetnikov,
\emph{Concise presentations of intuitionistic primitive recursive
arithmetic},
\url{https://github.com/VladimirReshetnikov/ProveIt/tree/25da87848ad4042c318af9bbc7f2e1c24614d6c3/Logic/PrimitiveRecursiveArithmetic/ConciseIPRA}.
The intended transfer endpoint is
\code{Lean/ConciseIPRA/NativeBiInterpretation.lean}.

\bibitem{MathlibAbelRuffini}
The mathlib community,
\code{Mathlib.FieldTheory.AbelRuffini} and
\code{Archive.Wiedijk100Theorems.AbelRuffini}, in the Lean mathematical
library.  ProveIt's semantic obstruction imports this checked development.

\bibitem{MathCompAbel}
The Mathematical Components community,
\emph{Abel: a formal proof of the Abel--Ruffini theorem},
\url{https://github.com/math-comp/Abel}.
ProveIt's Rocq/Coq development uses this independent formalization.

\end{thebibliography}

\subsection*{Repository pin and interpretation}

All repository paths and theorem names in this report refer to the following
ProveIt revision, inspected on 15 November 2026:
\begin{center}
\small\texttt{25da87848ad4042c318af9bbc7f2e1c24614d6c3}.
\end{center}
Later repository revisions may rename modules or strengthen theorems.  The
mathematical argument of this report is independent of those names; the pin
exists to make the Lean reuse plan reproducible.


\end{document}
